\section*{Chapitre \ref{ch:b3:deformation-mechanisms} --- Mécanismes de déformation et zones de cisaillement}

\begin{solution}{exo:b3:deformation-mechanisms:1}
Forstérite~: $1573/2163 = 0.73$~; glace~: $263/273 = 0.96$. La glace est bien plus proche
de son point de fusion et flue plus facilement~: les glaciers
s'écoulent à la surface de la Terre.
\end{solution}

\begin{solution}{exo:b3:deformation-mechanisms:2}
$\gamma = \tan\ang{70} = 2.75$ et $\tan\ang{80} = 5.67$. Près de \ang{90}, la
tangente varie très vite avec l'angle~: une erreur d'un degré sur $\psi =
\ang{84}$ modifie $\gamma$ d'un cinquième.
\end{solution}

\begin{solution}{exo:b3:deformation-mechanisms:3}
$\sigma = (10^{3.56}/10)^{1/1.26} = 363.1^{0.794} = \qty{108}{MPa}$.
\end{solution}

\begin{solution}{exo:b3:deformation-mechanisms:4}
$E/RT = 223000/(8.314\times773.15) = 34.69$, donc $\dot\varepsilon =
\num{1.1e-4}\times20^4\times e^{-34.69} = \num{1.1e-4}\times\num{1.6e5}\times
\num{8.6e-16} = \qty{1.5e-14}{s^{-1}}$.
\end{solution}

\begin{solution}{exo:b3:deformation-mechanisms:5}
$\eta = \sigma/(2\dot\varepsilon)$~: $\num{65.4e6}/\num{2e-14} =
\qty{3.3e21}{Pa.s}$ à \qty{400}{\celsius} et $\num{18.0e6}/\num{2e-14} =
\qty{9.0e20}{Pa.s}$ à \qty{500}{\celsius}~: la croûte moyenne à ces
températures est à peu près aussi visqueuse que le manteau sous les calottes glaciaires.
\end{solution}

\begin{solution}{exo:b3:deformation-mechanisms:6}
$n = \log(\num{1.6e-4}/\num{e-5})/\log2 = \log16/\log2 = 4$~: \omterm{def:b3:deformation-mechanisms:creep}{fluage par dislocations},
dont l'\omterm{def:b3:deformation-mechanisms:flow-law}{exposant de contrainte} est largement supérieur à 1.
\end{solution}

\begin{solution}{exo:b3:deformation-mechanisms:7}
$\ln10 = (E/R)(1/1000 - 1/1100)$, donc $E = 8.314\times2.3026/\num{9.091e-5} =
\qty{211}{kJ.mol^{-1}}$.
\end{solution}

\begin{solution}{exo:b3:deformation-mechanisms:8}
Nabarro--Herring ($d^{-2}$)~: 100 fois plus rapide~; Coble ($d^{-3}$)~: 1000 fois.
Le long des joints de grains, les atomes qui diffusent empruntent un chemin dont la section, $\delta
d$, qui diminue avec la taille du grain, s'ajoute à la distance plus courte et au
volume plus petit à déplacer.
\end{solution}

\begin{solution}{exo:b3:deformation-mechanisms:9}
$\sigma_b^{2.5} = \num{3.59e-6}\,d^{-3}/\num{6.54e-17}$~: pour $d = 1000$,
$\sigma_b = 54.9^{0.4} = \qty{5.0}{MPa}$~; pour $d = 100$, $\sigma_b =
\num{54900}^{0.4} = \qty{79}{MPa}$. La pente vaut $-3/2.5 = -1.2$~: un grain
dix fois plus fin augmente la contrainte d'un facteur $10^{1.2} = 16$.
\end{solution}

\begin{solution}{exo:b3:deformation-mechanisms:10}
Dans les deux chemins, l'écart de concentration moteur est le même et la
distance vaut $d$~; le volume à déplacer par unité de déformation est $d^3$. Le réseau
offre une section $\propto d^2$, donc la vitesse varie comme $d^2/(d\,d^3) =
d^{-2}$~; les joints de grains offrent $\propto\delta d$, donc $\delta d/(d\,d^3) = \delta
d^{-3}$. Le rapport des vitesses aux joints et dans le réseau est $(B_B\delta D_B)/(B_LD_Ld)$,
supérieur à 1 lorsque $d < (B_B/B_L)\,\delta D_B/D_L$.
\end{solution}

\begin{solution}{exo:b3:deformation-mechanisms:11}
$\eta \propto \dot\varepsilon^{(1-n)/n} = \dot\varepsilon^{-2.5/3.5}$~: dix fois
plus rapide divise la viscosité par $10^{0.714} = 5.2$. Une bande qui se déforme
par hasard plus vite devient plus faible, si bien qu'elle prend une plus grande part de la
déformation~: elle s'y concentre.
\end{solution}

\begin{solution}{exo:b3:deformation-mechanisms:12}
Avant~: $\dot\varepsilon_{\mathrm{dis}} = \num{6.54e-17}\times10^{3.5} =
\qty{2.07e-13}{s^{-1}}$ et $\dot\varepsilon_{\mathrm{dif}} = \num{3.59e-6}
\times10/3000^3 = \qty{1.3e-15}{s^{-1}}$~: \qty{2.08e-13}{s^{-1}}, \omterm{def:b3:deformation-mechanisms:creep}{fluage par
dislocations}. Après~: $\dot\varepsilon_{\mathrm{dif}} = \num{3.59e-6}\times10/30^3 =
\qty{1.33e-9}{s^{-1}}$, \omterm{def:b3:deformation-mechanisms:creep}{fluage par diffusion}, \num{6400} fois plus rapide sous la même
contrainte~: la recristallisation a considérablement adouci la zone.
\end{solution}

\begin{solution}{pb:b3:deformation-mechanisms:1}
\textbf{1.} Le filon passe de \ang{90} à \ang{6} par rapport à la zone~: $\psi =
\ang{84}$ et $\gamma = \tan\ang{84} = 9.51$.

\textbf{2.} $\gamma w = 9.51\times400 = \qty{3804}{m}$, soit environ \qty{3.8}{km}.

\textbf{3.} $S_1 = \sqrt{1 + 9.51^2/4} + 9.51/2 = 4.859 + 4.755 = 9.614$,
$S_2 = 1/S_1 = 0.104$, et $S_1/S_2 = 92.4$.

\textbf{4.} $\tan2\theta' = 2/9.51 = 0.210$, $\theta' = \ang{5.9}$~: la
foliation, proche du grand axe, se rapproche de plus en plus de la zone à mesure que $\gamma$
augmente.

\textbf{5.} Suivi vers le nord, le filon est entraîné vers l'est à l'intérieur de la zone et
en ressort plus à l'est~: le mur nord s'est déplacé vers l'est par rapport au mur sud.
Vu depuis le mur sud, le côté opposé s'est déplacé vers la droite~: la zone est
dextre.

\textbf{6.} Dans le plan contenant la linéation d'étirement, horizontale ici
pour une zone de décrochement, et perpendiculaire à la foliation verticale~: une
coupe horizontale. Sa foliation S s'incline dans le sens antihoraire par rapport à la zone, vers
l'est-nord-est, comme le grand axe de l'ellipse de déformation.

\textbf{7.} $\sigma = (10^{3.56}/20)^{1/1.26} = 181.5^{0.794} = \qty{62}{MPa}$.

\textbf{8.} $242.1^{0.794} = \qty{78}{MPa}$ pour \qty{15}{\micro m} et
$145.2^{0.794} = \qty{52}{MPa}$ pour \qty{25}{\micro m}.

\textbf{9.} Oui~: les expériences couvraient de 34 à \qty{268}{MPa}.

\textbf{10.} La zone s'est déformée à \qty{450}{\celsius}, bien en dessous des 800 à
\qty{1100}{\celsius} des expériences~; une relation qui ne dépend pas de la
température peut être extrapolée jusqu'à celle-ci, même s'il s'agit toujours
d'une extrapolation.

\textbf{11.} $D = 3631/100^{1.26} = 3631/331 = \qty{11}{\micro m}$.

\textbf{12.} $E/RT = 223000/(8.314\times723.15) = 37.09$ et $e^{-37.09} =
\num{7.79e-17}$.

\textbf{13.} $\dot\varepsilon = \num{1.1e-4}\times62^4\times\num{7.79e-17} =
\num{1.1e-4}\times\num{1.478e7}\times\num{7.79e-17} = \qty{1.27e-13}{s^{-1}}$.

\textbf{14.} $\dot\gamma = 2\dot\varepsilon = \qty{2.53e-13}{s^{-1}}$.

\textbf{15.} $\eta = \num{62e6}/(2\times\num{1.27e-13}) = \qty{2.4e20}{Pa.s}$,
proche des \qty{3e20}{Pa.s} du manteau supérieur sous le lac Ardo~: la zone
était aussi faible que le manteau.

\textbf{16.} $e^{-E/RT}$ est multiplié par 0.064 à \qty{400}{\celsius}
(soit seize fois plus lent) et par 11.0 à \qty{500}{\celsius}.

\textbf{17.} $(78/62)^4 = 2.5$ fois plus rapide.

\textbf{18.} $t = 9.51/\num{2.53e-13} = \qty{3.75e13}{s}$.

\textbf{19.} $\num{3.75e13}/\num{3.156e10} = \qty{1190}{ka}$.

\textbf{20.} $\qty{3804}{m}/\qty{1.19e6}{an} = \qty{3.2}{mm/an}$.

\textbf{21.} Dix fois plus lent, voire davantage, que les principales limites de plaques~: la zone
était un accident mineur, ou l'une de plusieurs structures se partageant le mouvement d'une limite.

\textbf{22.} $\dot\gamma = 0.02/(\num{3.156e7}\times400) = \qty{1.58e-12}{s^{-1}}$,
6.2 fois plus rapide~; la contrainte augmente de $6.2^{1/4} = 1.58$, pour atteindre \qty{98}{MPa}, et
les grains seraient de $3631/98^{1.26} = \qty{11}{\micro m}$.

\textbf{23.} Le temps varie comme $\sigma^{-4}$~: la seule plage de taille de grains donne
$(62/78)^4 = 0.40$ à $(62/52)^4 = 2.0$ fois l'estimation, et l'incertitude de
l'exposant $1.26 \pm 0.13$ ainsi que celle des constantes de la \omterm{def:b3:deformation-mechanisms:flow-law}{loi de fluage}
l'élargissent encore~: le résultat est un ordre de grandeur, non une date.

\textbf{24.} La \omterm{def:b3:deformation-mechanisms:shear-zone}{zone de cisaillement} de Varden a mis environ
\textbf{\qty[text-series-to-math=true]{1190}{ka}} pour accumuler sa déformation.
\end{solution}
